% Rozdział 5 — WERSJA „PREZENTACJA” (main-prezentacja.tex), 22.11.2026 (kris):
% same wyniki z parametrami pomiarów, bez interpretacji, z pustym
% podrozdziałem „Wnioski”. Pełna wersja: rozdzialy/05_badania.tex.
% Dane i rysunki te same co w emisja_tresc.tex, widma_tresc.tex
% i widma_materialy_tresc.tex — przy zmianach danych poprawić oba pliki.
\chapter{Badania eksperymentalne}
\label{ch:badania}

\renewcommand{\figdir}{../emisja-magnes-pomiary/figures}
\renewcommand{\datadir}{../emisja-magnes-pomiary/dane}
\providecommand{\sesjaM}{z~magnesem}
\providecommand{\sesjaB}{bez magnesu}
\pgfplotsset{/pgf/number format/use comma, /pgf/number format/1000 sep={}}

W~rozdziale przedstawiono wyniki pomiarów emisji polowej prototypu
z~magnesem ogniskującym i~bez niego oraz widma promieniowania X
zarejestrowane dla próbek o~różnym składzie. Stanowisko i~metody opisano
w~\rozdzref{ch:stanowisko}.

\section{Wpływ magnesu ogniskującego na emisję polową}
\label{sec:emisja-magnes}

\subsection{Przebieg pomiarów}

Wykonano sesje pomiarowe, w~których przy stałym napięciu $\UT$
zwiększano napięcie $\UB$ co \SI{100}{V} i~rejestrowano prąd emisji
$\Ie$ oraz zdjęcie ekranu luminescencyjnego (\rozdzref{ch:stanowisko}):
\begin{itemize}
  \item \textbf{sesja bez magnesu} -- $\UT=\SI{2000}{V}$,
        $\UB$ od \SI{500}{V} do \SI{1500}{V}, zdjęcia dla $\UB$ od
        \SI{700}{V} do \SI{1100}{V}; przed sesją źródło odgazowano
        przez wygrzanie w~temperaturze ok. \SI{200}{\celsius},
  \item \textbf{sesja z~magnesem} -- $\UT=\SI{1500}{V}$,
        $\UB$ od \SI{500}{V} do \SI{1700}{V}, zdjęcia dla $\UB$ od
        \SI{500}{V} do \SI{1600}{V}; bez ponownego odgazowania źródła,
  \item \textbf{seria z~magnesem przy $\UT=\SI{10}{kV}$} --
        $\UB$ od \SI{600}{V} do \SI{1400}{V}, bez zdjęć ekranu.
\end{itemize}

Dodatkowo wykonano serię~2 bez magnesu z~odczytami prądu od
\SI{1}{\micro A} do \SI{400}{\micro A}. Ciśnienie w~komorze nie
przekraczało \SI{1e-5}{\milli\bar}. Zmierzone prądy zestawiono
w~tab.~\ref{tab:emisja-iv}.

\begin{table}[htbp]
  \centering
  \caption{Prąd emisji w~funkcji napięcia $\UB$ w~sesjach pomiarowych}
  \label{tab:emisja-iv}
  \begin{tabular}{@{}r r r r@{}}
    \toprule
    $\UB$ [V] & $\Ie$ \sesjaB\ [µA] & $\Ie$ \sesjaM\ [µA] & $\Ie$ \sesjaM\ [µA] \\
    \midrule
     500 & 0    & 14 & -- \\
     600 & 0    & 16 & 4 \\
     700 & 0    & 14 & 18 \\
     800 & 2    & 28 & 36 \\
     900 & 3,1  & 45 & 44 \\
    1\,000 & 4    & 69 & 99 \\
    1\,100 & 6    & 118 & 127 \\
    1\,200 & 8    & 180 & 127 \\
    1\,300 & 9    & 190 & 126 \\
    1\,400 & 11,2 & 190 & 125 \\
    1\,500 & 22   & 193 & -- \\
    1\,600 & --   & 200 & -- \\
    1\,700 & --   & 200 & -- \\
    \midrule
    $\UT$ & \SI{2000}{V} & \SI{1500}{V} & \SI{10}{kV} \\
    \bottomrule
  \end{tabular}
\end{table}

\subsection{Charakterystyki prądowo-napięciowe}

Charakterystyki $\Ie(\UB)$ wszystkich sesji przedstawiono na rys.~\ref{fig:emisja-iv}.

\begin{figure}[htbp]
  \centering
  \begin{tikzpicture}
    \begin{semilogyaxis}[
      width=0.85\textwidth, height=7cm,
      xlabel={napięcie $\UB$ [V]}, ylabel={prąd emisji $\Ie$ [µA]},
      xmin=450, xmax=1750, ymin=1, ymax=500, xtick distance=200,
      grid=major, grid style={gray!25},
      legend pos=south east, legend cell align=left,
      every axis plot/.append style={thick, mark size=2pt},
      unbounded coords=jump,
    ]
      \addplot[red!60!black, mark=triangle*] table[x=U,y=I_uA,col sep=comma]{\datadir/iv_bez.csv};
      \addlegendentry{\sesjaB\ ($\UT=\SI{2000}{V}$)}
      \addplot[blue!70!black, mark=*] table[x=U,y=I_uA,col sep=comma]{\datadir/iv_mag.csv};
      \addlegendentry{\sesjaM\ ($\UT=\SI{1500}{V}$)}
      \addplot[black, mark=diamond*, mark size=2.6pt] table[x=U,y=I_uA,col sep=comma]{\datadir/iv_10kV.csv};
      \addlegendentry{\sesjaM\ ($\UT=\SI{10}{kV}$)}
    \end{semilogyaxis}
  \end{tikzpicture}
  \caption{Prąd emisji w~funkcji napięcia $\UB$: \sesjaB\ przy
    $\UT=\SI{2000}{V}$, \sesjaM\ przy $\UT=\SI{1500}{V}$
    i~$\UT=\SI{10}{kV}$
    (skala logarytmiczna; punkty z~$\Ie=0$ pominięto)}
  \label{fig:emisja-iv}
\end{figure}

\section{Obserwacja wiązki na ekranie luminescencyjnym}
\label{sec:emisja-ekran}

Wpływ magnesu na kształt wiązki oceniono na podstawie zdjęć ekranu
luminescencyjnego wykonanych w~sesjach opisanych w~\pref{sec:emisja-magnes}.

\subsection{Jasność ekranu}

Jasność świecenia ekranu $L$ wyznaczono ze zdjęć metodą opisaną
w~\pref{sec:analiza-zdjec}. Jej zależność od napięcia $\UB$ w~obu
sesjach przedstawiono na rys.~\ref{fig:emisja-I}, a~zależność od
odczytu prądu w~serii~2 na rys.~\ref{fig:emisja-IvsIe}.

\begin{figure}[htbp]
  \centering
  \begin{tikzpicture}
    \begin{axis}[
      width=0.8\textwidth, height=6.5cm,
      xlabel={napięcie $\UB$ [V]}, ylabel={jasność $L$ [j.u.]},
      xmin=450, xmax=1650, ymin=0, ymax=3.5, xtick distance=200,
      grid=major, grid style={gray!25},
      legend pos=north west, legend cell align=left,
      every axis plot/.append style={thick, mark size=2pt},
    ]
      \addplot[blue!70!black, mark=*] table[x=U,y=I,col sep=comma]{\datadir/mg.csv};
      \addlegendentry{\sesjaM\ ($\UT=\SI{1500}{V}$)}
      \addplot[red!60!black, mark=triangle*] table[x=U,y=I,col sep=comma]{\datadir/s1.csv};
      \addlegendentry{\sesjaB\ ($\UT=\SI{2000}{V}$)}
    \end{axis}
  \end{tikzpicture}
  \caption{Jasność świecenia ekranu w~funkcji napięcia $\UB$:
    \sesjaM\ przy $\UT=\SI{1500}{V}$, \sesjaB\ przy $\UT=\SI{2000}{V}$}
  \label{fig:emisja-I}
\end{figure}

\begin{figure}[htbp]
  \centering
  \begin{tikzpicture}
    \begin{loglogaxis}[
      width=0.65\textwidth, height=5.5cm,
      xtick={30,50,100,200,400}, xticklabels={30,50,100,200,400},
      log ticks with fixed point,
      xlabel={prąd emisji $\Ie$ [µA]},
      ylabel={jasność $L$ [j.u.]},
      grid=major, grid style={gray!25},
    ]
      \addplot[only marks, mark=square*, orange!85!black]
        table[x=etykieta, y=I, col sep=comma,
              restrict expr to domain={\thisrow{I}}{0.01:100}]{\datadir/s2.csv};
    \end{loglogaxis}
  \end{tikzpicture}
  \caption{Jasność świecenia ekranu w~funkcji odczytu prądu emisji,
    seria~2 \sesjaM}
  \label{fig:emisja-IvsIe}
\end{figure}

\subsection{Obraz plamki}

Zdjęcia ekranu z~obu sesji przy napięciach $\UB$ od \SI{800}{V} do
\SI{1100}{V} zestawiono na rys.~\ref{fig:emisja-sbs}, zdjęcia z~sesji
z~magnesem przy \SI{700}{V}, \SI{1400}{V} i~\SI{1600}{V} na
rys.~\ref{fig:emisja-mg-skrajne}, a~pełne kadry na
rys.~\ref{fig:emisja-pelne}. Parametry emisji i~plamki podano
w~tab.~\ref{tab:emisja}.

\begin{table}[htbp]
  \centering
  \caption{Parametry emisji i~plamki w~obu sesjach}
  \label{tab:emisja}
  \begin{tabular}{@{}lccc@{}}
    \toprule
    Parametr & \sesjaM & \sesjaB & seria 2 (\sesjaM) \\
    \midrule
    Odgazowanie źródła & nie & tak (\SI{200}{\celsius}) & tak (\SI{200}{\celsius}) \\
    Napięcie $\UT$ & \SI{1500}{V} & \SI{2000}{V} & \SI{2000}{V} \\
    Pierwsza widoczna plamka & \SI{600}{V} & \SI{800}{V} & \SI{700}{V} \\
    Prąd maksymalny & \SI{200}{\micro A} & \SI{22}{\micro A} & \SI{400}{\micro A} \\
    Maks. jasność $L$ [j.u.] & 3,19 & 1,67 & 2,89 \\
    Maks. pole plamki [\% kadru] & 3,98 & 2,07 & 3,91 \\
    \bottomrule
  \end{tabular}
\end{table}

\newcommand{\sbs}[2]{%
  \begin{subfigure}[t]{0.28\textwidth}\centering
    \includegraphics[width=\linewidth]{\figdir/#1}\caption{#2}
  \end{subfigure}}

\begin{figure}[p]
  \centering
  \makebox[0.28\textwidth]{\textbf{\sesjaM}}\hspace{1.5em}\makebox[0.28\textwidth]{\textbf{\sesjaB}}\par\medskip
  \sbs{mg_03.jpg}{\SI{800}{V}, \SI{28}{\micro A}}\hspace{1.5em}
  \sbs{s1_02.jpg}{\SI{800}{V}, \SI{2}{\micro A}}\par\medskip
  \sbs{mg_04.jpg}{\SI{900}{V}, \SI{45}{\micro A}}\hspace{1.5em}
  \sbs{s1_03.jpg}{\SI{900}{V}, \SI{3,1}{\micro A}}\par\medskip
  \sbs{mg_05.jpg}{\SI{1000}{V}, \SI{69}{\micro A}}\hspace{1.5em}
  \sbs{s1_04.jpg}{\SI{1000}{V}, \SI{4}{\micro A}}\par\medskip
  \sbs{mg_06.jpg}{\SI{1100}{V}, \SI{118}{\micro A}}\hspace{1.5em}
  \sbs{s1_05.jpg}{\SI{1100}{V}, \SI{6}{\micro A}}
  \caption{Zbliżenia okna ekranu \sesjaM\ ($\UT=\SI{1500}{V}$, lewa kolumna)
    i~\sesjaB\ ($\UT=\SI{2000}{V}$, prawa kolumna) przy tych samych
    napięciach $\UB$; podpisy: $\UB$, $\Ie$}
  \label{fig:emisja-sbs}
\end{figure}

\begin{figure}[htbp]
  \centering
  \begin{subfigure}[t]{0.28\textwidth}\centering
    \includegraphics[width=\linewidth]{\figdir/mg_02.jpg}\caption{\SI{700}{V}}\end{subfigure}\hspace{1.5em}
  \begin{subfigure}[t]{0.28\textwidth}\centering
    \includegraphics[width=\linewidth]{\figdir/mg_09.jpg}\caption{\SI{1400}{V}}\end{subfigure}\hspace{1.5em}
  \begin{subfigure}[t]{0.28\textwidth}\centering
    \includegraphics[width=\linewidth]{\figdir/mg_11.jpg}\caption{\SI{1600}{V}}\end{subfigure}
  \caption{Ekran w~sesji \sesjaM\ ($\UT=\SI{1500}{V}$) przy
    $\UB = \SI{700}{V}$, \SI{1400}{V} i~\SI{1600}{V}}
  \label{fig:emisja-mg-skrajne}
\end{figure}

\begin{figure}[htbp]
  \centering
  \begin{subfigure}[t]{0.32\textwidth}\centering
    \includegraphics[width=\linewidth]{\figdir/mg_pelny.jpg}\caption{\sesjaM, \SI{1400}{V}}\end{subfigure}\hfill
  \begin{subfigure}[t]{0.32\textwidth}\centering
    \includegraphics[width=\linewidth]{\figdir/s1_pelny.jpg}\caption{\sesjaB, \SI{1100}{V}}\end{subfigure}\hfill
  \begin{subfigure}[t]{0.32\textwidth}\centering
    \includegraphics[width=\linewidth]{\figdir/s2_pelny.jpg}\caption{\sesjaM, seria 2}\end{subfigure}
  \caption{Pełne kadry ekranu przy największym prądzie w~każdej serii}
  \label{fig:emisja-pelne}
\end{figure}

\subsection{Wnioski}

\vspace{6\baselineskip}

\section{Widma EDS próbek}
\label{sec:widma}

\subsection{Widmo próbki aluminiowej przy $\UT = \SI{4}{\kilo\volt}$}
\label{sec:widma-warunki}\label{sec:widma-linie}\label{sec:widma-stabilnosc}

Widmo zarejestrowano 12.06.2026 dla płytki aluminiowej, z~magnesem
ogniskującym na emiterze, przy prądzie emisji $\Ie \approx
\SI{100}{\micro\ampere}$ i~ciśnieniu nie większym niż
\SI{1e-5}{\milli\bar}, spektrometrem opisanym
w~\pref{sec:rejestracja-widm}. Parametry pomiaru zestawiono
w~\tabref{tab:widma-parametry}, widmo przedstawiono na
\rysref{fig:widmo-4kV}, zidentyfikowane linie w~\tabref{tab:widma-linie},
a~przebieg częstości zliczeń w~czasie pomiaru na \rysref{fig:widma-stabilnosc}.

\begin{table}[htbp]\centering
  \caption{Parametry rejestracji widma EDS}\label{tab:widma-parametry}
  \begin{tabular}{@{}ll@{}}\toprule
    Parametr & Wartość \\\midrule
    Napięcie przyspieszające $\UT$ & \SI{4000}{\volt} \\
    Napięcie ekstrakcyjne $\UB$ & \SI{1000}{\volt} \\
    Prąd emisji $\Ie$ & ok.~\SI{100}{\micro\ampere} \\
    Czas kształtowania (\angielski{peaking time}) & \SI{0.5}{\micro\second} \\
    Liczba kanałów, szerokość kanału & \num{8192}, \SI{10}{\electronvolt} \\
    Czas rzeczywisty / czas życia & \SI{210.7}{\second} / \SI{100.1}{\second} \\
    Częstość zliczeń wejściowa / wyjściowa & ok.~475~kcps / ok.~225~kcps \\
    Czas martwy & \SIrange{52}{54}{\percent} \\\bottomrule
  \end{tabular}
\end{table}

\rysunek[0.95\linewidth]{widma-eds/widmo_eds_4kV.pdf}{Widmo EDS płytki
  aluminiowej przy $\UT=\SI{4}{\kilo\volt}$ i~$\UB=\SI{1}{\kilo\volt}$,
  czas życia \SI{100}{\second}}{fig:widmo-4kV}

\begin{table}[htbp]\centering
  \caption{Linie zidentyfikowane w~widmie z~\rysref{fig:widmo-4kV}}\label{tab:widma-linie}
  \begin{tabular}{@{}S[table-format=1.2]lS[table-format=1.3]@{}}\toprule
    {$E$ piku [\si{\kilo\electronvolt}]} & Linia & {$E$ tabel. [\si{\kilo\electronvolt}]} \\\midrule
    1.26 & Mg~\Ka & 1.254 \\
    1.48 & Al~\Ka & 1.487 \\
    1.74 & Si~\Ka & 1.740 \\
    2.98 & Al~\Ka{} + Al~\Ka & 2.974 \\
    3.22 & Al~\Ka{} + Si~\Ka & 3.227 \\
    3.46 & Si~\Ka{} + Si~\Ka & 3.480 \\\bottomrule
  \end{tabular}
\end{table}

\rysunek[0.9\linewidth]{widma-eds/stabilnosc_zliczen.pdf}{Częstość zliczeń
  wejściowa (ICR), wyjściowa (OCR) i~średnia częstość w~kanale Al~\Ka{}
  w~czasie rejestracji widma przy $\UT=\SI{4}{\kilo\volt}$
}{fig:widma-stabilnosc}

\subsection{Widma próbek o~różnym składzie przy $\UT = \SI{10}{\kilo\volt}$}
\label{sec:widma-probki}\label{sec:widma-tlen}

Widma zarejestrowano 13.11.2026 przy $\UT=\SI{10}{\kilo\volt}$ dla
czterech próbek: A~(miedź o~nieznanym składzie), B~(ołów), C~(złoto,
chrom i~miedź) i~D~(aluminium). Dla próbek A--C wykonano po 6--10 serii
o~czasie życia \SIrange{40}{250}{\second}, które zsumowano; dla
próbki~D wykorzystano serię o~czasie martwym \SI{20}{\percent}. Skalę
energii skorygowano liniowo na podstawie położeń linii Al~\Ka,
Pb~M$\alpha$, Cr~\Ka, Fe~\Ka{} i~Cu~\Ka. Widma przedstawiono na
\rysref{fig:widma-probki}, a~zidentyfikowane linie w~\tabref{tab:widma-probki}.

\rysunek[0.95\linewidth]{widma-eds/widma_materialy.pdf}{Widma EDS próbek A, B i~C
  przy $\UT=\SI{10}{\kilo\volt}$ (13.11.2026, skala logarytmiczna, zliczenia
  na sekundę czasu życia)}{fig:widma-probki}

\begin{table}[htbp]\centering\small
  \caption{Linie zidentyfikowane w~widmach z~13.11.2026}\label{tab:widma-probki}
  \begin{tabular}{@{}lll@{}}\toprule
    Próbka & Pierwiastek & Linie i~energie po korekcji [\si{\kilo\electronvolt}] \\\midrule
    A & Cu       & L$\alpha$ 0,92; \Ka{} 8,04; \Kb{} 8,90 \\
      & Te       & L$\alpha$ 3,76; L$\beta_1$ 4,04; L$\beta_2$ 4,30; L$\gamma_1$ 4,57 \\
      & P        & \Ka{} 2,01 \\
      & S, Cl, K & \Ka{} 2,30; 2,62; 3,32 \\
      & Si, Al   & \Ka{} 1,73; 1,48 \\\midrule
    B & Pb       & M$\zeta$ 1,83; M$\alpha$ 2,35; M$\gamma$ 2,65 \\\midrule
    C & Au       & M$\zeta$ 1,66; M$\alpha\beta$ 2,13; M$\gamma$ 2,41 \\
      & Cr       & \Ka{} 5,42; \Kb{} 5,95 \\
      & Cu       & L$\alpha$ 0,93; \Ka{} 8,04 \\\midrule
    D & Al       & \Ka{} 1,48 \\
      & Mn, Fe   & \Ka{} 5,91; 6,40 \\
      & --       & pik sumacyjny 2,98 \\\bottomrule
  \end{tabular}
\end{table}
